% BHU PYQ -- Professional Exam Template
\documentclass[11pt,a4paper]{article}

\usepackage[a4paper,top=2.5cm,bottom=2.5cm,left=2.8cm,right=2.8cm,headheight=14pt]{geometry}
\usepackage{amsmath,amssymb}
\usepackage[T1]{fontenc}
\usepackage[utf8]{inputenc}
\usepackage{lmodern}
\usepackage{microtype}
\usepackage{fancyhdr}
\usepackage{enumitem}
\usepackage{listings}
\usepackage{hyperref}
\usepackage{lastpage}
\usepackage{parskip}

\hypersetup{colorlinks=true,urlcolor=black,linkcolor=black,
  pdftitle={CS-109: Data Communication -- BHU PYQ}}

\pagestyle{fancy}
\fancyhf{}
\renewcommand{\headrulewidth}{0.4pt}
\renewcommand{\footrulewidth}{0.4pt}
\fancyhead[L]{\small\textit{Banaras Hindu University}}
\fancyhead[R]{\small\textit{CS-109: Data Communication}}
\fancyfoot[C]{\small Page \thepage\ of \pageref{LastPage}}

\lstset{
  basicstyle=\small\ttfamily,
  frame=single,
  breaklines=true,
  columns=flexible,
  keepspaces=true
}

\newlist{parts}{enumerate}{2}
\setlist[parts,1]{label=(\alph*),leftmargin=*,itemsep=4pt,topsep=3pt}
\setlist[parts,2]{label=(\roman*),leftmargin=*,itemsep=2pt,topsep=2pt}
\newcommand{\pts}[1]{\hfill\textit{[#1]}}

\begin{document}
\begin{center}
  {\large\bfseries BANARAS HINDU UNIVERSITY}\\[2pt]
  {\normalsize B.Sc. (Hons.) Semester VI Examination, 2022-23}\\[6pt]
  \rule{0.8\linewidth}{0.5pt}\\[6pt]
  {\large\bfseries Paper: CS-109 --- Data Communication}\\[4pt]
  {\normalsize Time: Three hours \hfill Maximum Marks: 70}
\end{center}

\rule{\linewidth}{0.4pt}

\smallskip
\noindent\textit{Instructions: ANSWER ANY FIVE QUESTIONS, ALL QUESTIONS CARRY EQUAL MARKS.}

\bigskip

CS - 109 : Data Communication

\medskip\noindent\textbf{Question 1.} Write a short note on any four out of six techniques/protocols: 3.5x4=14

\medskip\noindent\textbf{Question 2.} Amplitude Modulation(AM) ,

\medskip\noindent\textbf{Question 5.} Frequency-Division Multiplexing (FDM) : : ‘
\begin{parts}
  \item Delta Modulation (DM)
  \item Cyclic Redundancy Check (CRC) .
  \item High Level Data Link Control (HDLC) .
  \item Direct Sequence Spread Spectrum (DSSS)
\end{parts}

\medskip\noindent\textbf{Question 2.}
\_ a) What is wireless (unguided) transmission media? What are the three major types of wireless
- transmission media, describe each in detail. 7 Cee 6 “
\begin{parts}
  \item What are circuit-switched and packet-switched networks? Discuss differences among them. 8
\end{parts}

\medskip\noindent\textbf{Question 3.} .
\begin{parts}
  \item What is the result of scrambling the sequence 11100000000000 using one of the following
  scrambling techniques? Assume that the last non-zero signal level has been positive. 6
  \item B8ZS (ii) HDB3 (The number of nonzero pulses is odd after the last substitution)
  \item What is the purpose of Pulse Code Modulation (PCM) technique? Describe the components of a
  PCM encoder with an example. c 8:
\end{parts}

\medskip\noindent\textbf{Question 4.} . .
\begin{parts}
  \item Describe the following mechanisms for modulating digital data into an analog signal with examples:
  Amplitude shift keying (ASK), Frequency shift keying (FSK), Phase shift keying (PSK), and
  Quadrature amplitude modulation (QAM). 6
  \item Station 4 needs to send a message consisting of 9 packets to Station B using asliding window
  (window size 3) and go-back-n error control strategy. All packets are ready and immediately
  available for transmission. If every Sth packet that A transmits gets lost (but no acks from B ever get
  lost), then what is the total number of packets that A will transmit for sending the message to B?
  Explain through diagram. 8
\end{parts}

\medskip\noindent\textbf{Question 5.}
\begin{parts}
  \item Discuss Frequency Hopping Spread Spectrum (FHSS) with a suitable diagram. 6
  \item A sender needs to send the four data items given in hexadecimal: 3456, ABCC, 02BC, and EEEE.
  Answer the following questions: 8
  : i. Calculate the checksum at the sender site.
  ii, Calculate the checksum at the receiver site if the second data item is changed to ABCE.
  \textasciitilde{} iil. ‘Find the checksum at the receiver site if the second data item is changed to ABCE and the
  third data item is changed to 02BA.
\end{parts}

\medskip\noindent\textbf{Question 6.}
\begin{parts}
  \item What is data link control? Describe the key components of data link control. 6
  \item Why we need flow and error control? Discuss the sfop-and-wait flow control technique with an
  example. 8
  Seek END eee eRe |
\end{parts}

\vfill
\rule{\linewidth}{0.4pt}
\begin{center}\small
\href{https://bhu-exam.vercel.app/}{Click here to visit our website}\\[4pt]\href{https://me-aryan.vercel.app/}{Click here to visit our contributor}\\[4pt]\href{https://quashan-me.vercel.app/}{Click here to visit our contributor}
\end{center}
\end{document}